% ============================================================
% ch31.tex —— 第 31 章 平行线的葬礼：射影平面
% ============================================================
\chapter{平行线的葬礼：射影平面}

第 30 章我们把"无穷远点"当救兵请来补贝祖定理的洞。这一章走进这片新大陆本身：它不玄——画家用了五百年；它不空——五个小节的代数能把你带到无穷远点现场，亲手结清前两章的每笔悬案。

\section{消失点：画家先于数学家五百年}

文艺复兴佛罗伦萨，布鲁内莱斯基（约 1413 年，那次著名的镜面透视演示）与阿尔贝蒂（1435 年《论绘画》）面对一个技术难题：怎么把三维世界\textbf{忠实}画到二维画布上？他们的答案成了西方绘画的分水岭——\textbf{透视法}。核心规律你天天见：

\textbf{平行的铁轨、走廊边线、林荫道两侧，在画面里汇聚于同一个"消失点"。}方向不同的平行线族，各有各的消失点（走廊的进深线灭于正前方的点，天花板横梁灭于别处）；所有"水平方向"的消失点连成一条线，叫\textbf{地平线}。画家的操作手册：先定地平线，再定各方向消失点，一切透视线奔各自的点而去——三百年后数学家一看：\textbf{这不就是一套几何公理吗}。

数学家的扶正工作只有两条：给每个方向的消失点起名"无穷远点"；再把所有无穷远点（地平线的抽象）合成"无穷远直线"。普通平面 $+$ 无穷远直线 $=$ \textbf{射影平面}。在这里，\textbf{任何两条不同直线恰好交于一点}——相交的老地方见，平行的去无穷远握手。欧氏几何里"平行线不相交"的公设被温和地废除：\textbf{平行只是"交点在无穷远"的别名}。

还有一条更深的礼物：在射影平面里，\textbf{点与直线地位完全平等}（"两点定一线"的对偶"两线定一点"也永远成立——因为平行线也有了交点）。几何学家由此开发出"对偶原理"：把定理里"点"与"线"互换，新定理自动成立（帕斯卡定理与布里昂雄定理就是这样一对双胞胎）。这份优雅是欧氏平面给不了的。

\section{选读小节：齐次坐标——给无穷远点发身份证}

（本小节五分钟，用完即走；跳过不影响正文。）

无穷远点要有用，得能进方程、能参与计算。办法：把坐标从两个数升级成\textbf{三个数} $[X : Y : Z]$（$Z$ 不全为零，即三数不同时为零），规则只有一条：
\begin{jielun}
三个数同乘一个非零常数，视为\textbf{同一点}：$[1:2:1] = [2:4:2] = [0.5:1:0.5]$。
\end{jielun}
（好比名片放大一号还是同一张名片。）普通点与老坐标的对接：$(x, y) = [x : y : 1]$（即 $Z = 1$ 那一层）。而 $Z = 0$ 的点（如 $[1:2:0]$、$[1:0:0]$）\textbf{写不回} $(x,y)$——\textbf{它们正是无穷远点}：$[X:Y:0]$ 读作"斜率 $\tfrac YX$ 的那个方向"（$X \ne 0$ 时），$[0:1:0]$ 是竖直方向。

把曲线搬进来只需一步：\textbf{每项补齐到同样的总次数}（乘上够多的 $Z$）。例：
\begin{gather*}
y = 1 \;\to\; Y = Z; \qquad y = 0 \;\to\; Y = 0; \\
xy = 1 \;\to\; XY = Z^2; \qquad x^2 + y^2 = 1 \;\to\; X^2 + Y^2 = Z^2 .
\end{gather*}
（检验次数：$XY$ 总次数 $2$，补 $Z^2$ 也是 $2$，齐了。）现在\textbf{亲手结案前两章的三笔账}：

\paragraph{账一（平行线，第 30 章实验三）。}$Y = 0$ 与 $Y - Z = 0$ 联立：相减得 $Z = 0$，代回 $Y = 0$，解出
\[
[X : Y : Z] = [1 : 0 : 0] \qquad (\text{水平方向的无穷远点}) .
\]
一票，不多不少，$1\times1 = 1$ \checkmark。

\paragraph{账二（渐近线，第 29 章悬案）。}$Y = 0$ 与双曲线 $XY = Z^2$ 联立：$Y = 0$ 代入得
\[
0 = Z^2 \;\Longrightarrow\; Z = 0\ \text{（\textbf{二重解}）},\qquad X\ \text{任意（取 }1\text{）}.
\]
交点 $[1:0:0]$，\textbf{重数二}——$2 = 1\times2$ \checkmark。原来横轴与双曲线在无穷远处"平方级贴合"——渐近的意义在代数上原来是"二阶接触"。悬案结清。

\paragraph{账三（抛物线与无穷远，第 30 章预告）。}$y = x^2$ 补齐：$YZ = X^2$。与无穷远直线 $Z = 0$ 联立：$X^2 = 0$，得 $[0:1:0]$（竖直方向），\textbf{重数二}——抛物线在无穷远处"自切"（它就是第 31 章选读小节的"在无穷远碰了一下自己"）。双曲线 $XY = Z^2$ 与 $Z = 0$ 联立：$XY = 0$，两个解 $[1:0:0]$ 与 $[0:1:0]$——\textbf{两个不同的无穷远点}，两臂各奔一个方向。

\section{圆的共同秘密}

齐次坐标立刻爆出一个冷知识。圆 $x^2 + y^2 = r^2$ 补齐成 $X^2 + Y^2 = r^2Z^2$，让它与无穷远直线 $Z = 0$ 联立：
\[
X^2 + Y^2 = 0 \quad\Longrightarrow\quad \frac{Y}{X} = \pm\mathrm{i} \quad\Longrightarrow\quad [1 : \mathrm{i} : 0] \text{ 与 } [1 : -\mathrm{i} : 0] .
\]
\textbf{每一个圆都穿过同样两个复数无穷远点}（外号\textbf{循环点}，circular points）——不管半径多大、圆心在哪。这解释了一桩几何事实：为什么"任意不共线三点定一圆"（圆的射影自由度：一般圆 $x^2+y^2+Dx+Ey+F=0$ 是 $3$ 个自由参数，三点三方程，恰好钉死），也解释了"所有圆彼此相似"的深层原因——它们共享两个藏在无穷远的顶点。贝祖记账也顺了：两个圆都是二次，$2\times2 = 4$ 票；循环点占两票（每对圆公共），剩两票给实数世界至多两个交点——\textbf{你从没在图上看见的那两张票，代数一直替你保管着}。

\section{圆锥曲线大团圆}

同一台机器，把三种圆锥曲线的辈分关系一次说破。看每条二次曲线与\textbf{无穷远直线} $Z=0$ 的交点数（用第 29 章选读的判别式 $B^2 - 4AC$ 的符号裁决）：

\begin{center}
\begin{tabular}{lccc}
\toprule
 & 与无穷远直线的关系 & 实数形象 & $B^2-4AC$ \\
\midrule
椭圆 & 两个\textbf{复}无穷远点（循环点，不真露面） & 闭口，永远有限 & $<0$ \\
抛物线 & \textbf{一个}无穷远点，重数二（相切） & 一端开口，"贴着"无穷 & $=0$ \\
双曲线 & \textbf{两个实}无穷远点（穿出去） & 两臂各奔一个方向 & $>0$ \\
\bottomrule
\end{tabular}
\end{center}

验算抛物线：$y = x^2$ 补齐 $YZ = X^2$，$Z = 0$ 代入得 $X^2 = 0$——单解重二（切）\checkmark。验算双曲线：$XY = Z^2$，$Z=0$ 给 $XY = 0$ 两解 \checkmark。验算圆：上面刚做，两复点 \checkmark。于是三种曲线在射影世界的关系一目了然：

\begin{jielun}
\textbf{抛物线是在无穷远处碰了一下自己的椭圆；双曲线是穿出无穷远的椭圆。}
\end{jielun}

三类曲线本是同根生（都是"二次曲线"），差别只在"与无穷远的关系"：躲着（椭圆）、贴着（抛物线）、穿过（双曲线）。阿波罗尼奥斯用切圆锥的手艺分门别类（第 29 章），齐次坐标一行代数统一三国——\textbf{升维一次，天下大同}，这是射影几何的招牌打法，也是本书"发明新工具简化老问题"的最后一次演出。

\begin{dongshou}
\begin{itemize}
  \yisi 把下列曲线补齐成齐次形式：$x^2 + y^2 = 25$；$y^2 = x$；$xy = 4$；$y = 2x + 3$。（答案：$X^2+Y^2 = 25Z^2$；$Y^2 = XZ$；$XY = 4Z^2$；$Y = 2X + 3Z$。注意直线是一次齐次。）
  \yier 求下列曲线的无穷远点并判断与无穷远直线的接触方式：$y^2 = x^3$（补齐 $Y^2Z = X^3$：$Z=0$ 给 $X^3 = 0$，单点 $[0:1:0]$ \textbf{重数三}——三次曲线与无穷远的"三阶接触"，这条曲线在无穷远处有个"尖点"，它是第 32 章椭圆曲线\textbf{反面教材}：三次项"重根"，群律失灵）；$x^2 - y^2 = 1$（$X^2 - Y^2 = Z^2$，无穷远给 $X^2 = Y^2$，两点 $[1:\pm1:0]$，双曲线两臂方向 \checkmark）。
  \yier 用"三点定一圆"的射影解释核对一个特例：为什么过\textbf{共线}三点没有圆？（三实点若共线，加上两个循环点共五点；"五点定一条二次曲线"（第 30 章烧脑时刻）在这里给出的二次曲线退化为"那条直线 $+$ 无穷远直线"——不是真的圆。三线共点破坏了"一般位置"。）
\end{itemize}
\end{dongshou}

\begin{shaonao}
\textbf{射影平面是单侧的。}一条纸带一端拧 $180^\circ$ 再首尾相接，得\textbf{莫比乌斯带}——蚂蚁沿带面爬一圈，回到起点时已在"另一面"（其实它从头到尾只有一面）。把莫比乌斯带的边缘（那条唯一的边）粘合成一个圆，得到的封闭曲面就是射影平面——\textbf{它没有内外两面}。齐次坐标早就暗示了这一点：$[X:Y:Z]$ 与 $[-X:-Y:-Z]$ 是同一点，即"方向"与"反方向"在无穷远被\textbf{粘在了一起}（水平向左与向右共享 $[1:0:0]$）。请做两件事：

(1) 用齐次坐标验证 $[2:0:0] = [1:0:0] = [-1:0:0]$（同乘 $-\tfrac12$）——"向左"与"向右"的无穷远点确实是同一个。

(2) 在莫比乌斯带上画一条"中线"，观察它两次回到起点——对应射影平面上"沿一条直线走无穷远再回来，方向翻转"。

顺带一个哲学彩蛋：射影平面在三维空间里造不出来（粘合过程必自相交），它只在四维里"住得下"——数学对象的合法性从不取决于"能不能捏出来"，\textbf{公理自洽即存在}（第 28 章的方法论，最后一次回响）。
\end{shaonao}

\begin{chengshi}
齐次坐标的完整理论（射影变换群、对偶原理的严格版、交比不变量）本章只取了"发身份证"一个功能。莫比乌斯带粘圆盘的构造在三维里必自交（严格证明要相交数理论）；"直线 $+$ 无穷远直线 $=$ 退化二次曲线"这类退化情形在"五点定二次曲线"里需要"一般位置"条款保护。循环点依赖复结构：实射影平面里圆与无穷远直线的交点账要复化后才平——这正是"复射影平面"成为代数几何标准舞台的原因。
\end{chengshi}

射影世界就绪，无穷远的户口本办齐。终章舞台搭好：一条三次曲线，即将同时拴住费马大定理、椭圆的周长与你的手机密码——而它的全部魔法，来自本章刚刚发证的"无穷远点 $O$"。
